\chapter{Lars Hörmander}\label{ch1962larshormander:lars-huxf6rmander-fields-medal-1962-a-comprehensive-report}

\begin{center}\rule{0.5\linewidth}{0.5pt}\end{center}

\textbf{Prepared:} July 2026
\textbf{Subject:} Lars Valter Hörmander (b. January 24, 1931, Mjällby, Sweden; d. November 2012, Lund, Sweden)
\textbf{Position:} Professor Emeritus, Lund University
\textbf{Doctoral Advisor:} Lars Gårding (Lund University)
\textbf{Fields Medal Citation:} \emph{``Worked in partial differential equations.
Specifically, contributed to the general theory of linear differential operators.
The questions go back to one of Hilbert's problems at the 1900 congress.''}

\begin{center}\rule{0.5\linewidth}{0.5pt}\end{center}

\begin{prereqs} Real analysis (Lebesgue measure, $L^p$ spaces), Fourier
analysis basics, basic functional analysis. \end{prereqs}

\begin{objectives} State the H\"{o}rmander condition for hypoellipticity; explain
the pseudodifferential operator framework; describe the $L^2$ theory of the
$\overline{\partial}$ operator. \end{objectives}

\section{Problem History}

The general theory of linear PDEs was one of Hilbert's 23 problems from the 1900
Paris Congress. Hilbert asked for a systematic theory of linear operators
$P(x, D) = \sum_{|\alpha| \leq m} a_\alpha(x) D^\alpha$ that would parallel the
well-developed theory of linear algebra.

A fundamental question was \textbf{hypoellipticity}: for which operators $P(D)$
does $P(D)u \in C^\infty$ imply $u \in C^\infty$? The Laplace operator is
hypoelliptic by elliptic regularity, the heat operator by parabolic theory, but
the wave operator is not.

Petrowsky studied the Cauchy problem for higher-order equations, identifying
conditions on the symbol $P(\xi)$ guaranteeing well-posedness. But the
relationship between symbol and regularity was not fully understood.

The \textbf{Cauchy problem}---solving $Pu = f$ with prescribed initial data---was
understood for hyperbolic equations but problematic for non-analytic data. Local
solvability was thought to always hold for linear operators.

In 1957, Hans Lewy shattered this belief with a smooth linear operator in
$\mathbb{R}^3$ having \emph{no local solutions anywhere}. The operator
$Lu = \frac{\partial u}{\partial \bar{z}} - i\bar{z}\frac{\partial u}{\partial t}$
has no distributional solution in any open set for most $f \in C^\infty$.

The problem of \textbf{several complex variables} needed new tools. The
$\bar{\partial}$-equation lacked the $L^2$ estimates needed for a
functional-analytic approach.

The \textbf{spectral theory} of elliptic operators was another frontier. Weyl's
law $N(\lambda) \sim C \cdot \text{Vol}(M) \lambda^{n/2}$ needed optimal remainder
estimates connecting analysis and geometry.

\section{How Was It Solved}

Hörmander's thesis \cite{hormander1963} (1955, \emph{Acta Mathematica}) gave a \textbf{complete
characterization of hypoellipticity} for constant-coefficient operators:
$P(D)$ is hypoelliptic if and only if:

\[\lim_{|\xi| \to \infty} \frac{|\partial^\alpha P(\xi)|}{|P(\xi)|^{1/|\alpha|}} = 0
\quad \text{for all } \alpha \neq 0\]

This means $P$ does not oscillate too rapidly relative to its magnitude. Elliptic
operators satisfy $|P(\xi)| \sim |\xi|^m$ and are automatically hypoelliptic.

For the $\bar{\partial}$-equation, Hörmander proved his \textbf{$L^2$ estimate
theorem} (1965). On pseudoconvex $\Omega \subset \mathbb{C}^n$ with plurisubharmonic
weight $\phi$: for $\bar{\partial}$-closed $(0,q)$-form $f$ with
$\int_\Omega |f|^2 e^{-\phi} dV < \infty$, there exists $u$ with $\bar{\partial}u = f$
and:

\[\int_\Omega |u|^2 e^{-\phi}(1+|z|^2)^{-2} dV \leq \int_\Omega |f|^2 e^{-\phi} dV\]

The proof uses the \textbf{Bochner--Kodaira--Nakano identity}:
$\Box = \nabla^*\nabla + [i\Theta(E), \Lambda]$.

The \textbf{Hörmander condition} (1967): if $L = X_1^2 + \cdots + X_k^2$ and the
Lie algebra generated by $X_1, \ldots, X_k$ spans $T_xM$ at every point, then $L$
is hypoelliptic. For the Heisenberg sub-Laplacian, $[X,Y] = \partial_t$ provides
the spanning.

The \textbf{$\Psi$DO calculus} extends differential operators via:

\[(Au)(x) = \frac{1}{(2\pi)^n}\int e^{ix\cdot\xi} a(x,\xi) \hat{u}(\xi) d\xi\]

with symbol classes $S^m_{\rho,\delta}$ and composition, adjoints, ellipticity, and
the sharp Gårding inequality.

\textbf{Fourier integral operators} (FIOs) generalize $\Psi$DOs for hyperbolic
equations. With Duistermaat, Hörmander proved \textbf{propagation of
singularities}: $\text{WF}(u)$ propagates along bicharacteristics.

Hörmander explained Lewy's phenomenon geometrically: solvability requires the
\textbf{Levi condition} $[X,Y] \in \text{span}(X,Y)$.

He proved the \textbf{sharp Weyl law}:

\[N(\lambda) = \frac{\omega_n}{(2\pi)^n}\text{Vol}(M)\lambda^{n/2} + O(\lambda^{(n-1)/2})\]

using FIOs and stationary phase.

\section{Impact of This Discovery}

Hörmander's hypoellipticity characterization answered one of Hilbert's problems
and established the modern framework for classifying PDEs.

The $L^2$ theory of $\bar{\partial}$ revolutionized several complex variables,
solving the Levi problem and providing Bergman kernel estimates.

The $\Psi$DO and FIO calculi became the language of modern PDE theory, essential
for the Atiyah--Singer index theorem and spectral geometry.

The Hörmander condition is fundamental in sub-Riemannian geometry and control
theory.

Propagation of singularities is the foundation of microlocal analysis, with
applications in seismic imaging and tomography.

Hörmander's four-volume treatise \emph{The Analysis of Linear Partial Differential
Operators} (1983--1985) is the definitive reference.

\section{Modern Views and Improvements}

The $\Psi$DO and FIO calculi have been refined for nonlinear problems by
Sjöstrand and Métivier.

The Hörmander condition has been generalized to sub-Laplacians on sub-Riemannian
manifolds with heat kernel asymptotics.

The Atiyah--Singer index theorem, using $\Psi$DOs and FIOs, has been extended to
families, manifolds with boundary, and foliations.

$L^2$ methods for $\bar{\partial}$ extend to CR geometry and the
$\bar{\partial}_b$-complex.

Active areas include resolvent estimates, Strichartz estimates, and Carleman
estimates for inverse problems.

\begin{center}\rule{0.5\linewidth}{0.5pt}\end{center}

\section*{Summary}
\addcontentsline{toc}{section}{Summary}

Hörmander developed the modern $L^2$ theory of PDEs, introducing pseudodifferential
operators and Fourier integral operators as the fundamental tools for analyzing
linear differential equations. His characterization of hypoellipticity and his work
on the $\overline{\partial}$-operator established the functional-analytic
foundations of modern PDE theory.
